% v01, EVB, 

%%%%%%%%%%%%%%%%%%%%%%%

\documentclass[a4paper]{article}
\usepackage[normalem]{ulem}
\usepackage{titlesec}
\setcounter{secnumdepth}{4}
\usepackage{comment}
\usepackage[english]{babel}
\usepackage[T1]{fontenc}
\usepackage{amssymb}
\usepackage{amsthm}
\usepackage[a4paper,top=2.5cm,bottom=2.5cm,left=2cm,right=2cm,marginparwidth=1.75cm]{geometry}
\usepackage{enumerate}
\usepackage{wrapfig}
\usepackage{enumerate}
\usepackage{enumitem}
\def\R{{\cal R}}
\usepackage[dvipsnames]{xcolor}

\newcounter{acounter}

\newcommand{\EV}[1]{\textcolor{red}{$\langle${\slshape{\bfseries Ev:} #1}$\rangle$}}
\newcommand{\evcolor}[1]{\textcolor{red}{#1}}
\newcommand{\map}[1]{\textcolor{blue}{#1}}
\newcommand{\joecolor}[1]{\textcolor{ForestGreen}{#1}}
\newcommand{\estrike}[1]{\textcolor{red}{\sout{#1}}}


\newcommand{\diag}{\operatorname{diag}}
\newcommand{\sign}{\operatorname{sign}}
\newcommand{\Ker}{\operatorname{Ker}}
\newcommand{\rank}{\operatorname{rank}}
\newcommand{\range}{\operatorname{Range}}
%\newcommand{\dg}{\operatorname{dg}}
\newcommand{\dg}{\operatorname{diag}}
\newcommand{\T}{\operatorname{\,T}}
\newcommand{\reals}{\mathbb{R}}
\newcommand{\NA}{\operatorname{na}}
\newcommand{\A}{\operatorname{a}}

\usepackage[export]{adjustbox}
\usepackage{subfig}
\usepackage{amsmath}
\usepackage{graphicx}
\usepackage[colorinlistoftodos]{todonotes}
\usepackage[colorlinks=true, allcolors=blue,hypertexnames=false]{hyperref}
\usepackage{float}

\newcommand{\SubItem}[1]{
    {\setlength\itemindent{15pt} \item[-] #1}
}

\usepackage{amsmath,amsfonts,amssymb,amsthm}
\usepackage{mathtools}
\usepackage{commath}
\let\oldnorm\norm  
\let\norm\undefined 
\DeclarePairedDelimiter\norm{\lVert}{\rVert}



\title{Transitions between mutualism, parasitism, predator-prey, and competition in an endosymbiont-host system}
\author{Essteban Vargas in collaboration with Michael Lynch and Adrian Gonzalez}
\newtheorem{theorem}{Theorem}
\newtheorem{definition}{Definition}
\newtheorem{proposition}{Proposition}
\newtheorem*{proposition*}{Proposition}

\newcommand{\joe}[1]{\textcolor{ForestGreen}{$\langle${\slshape{\bfseries Joe:} #1}$\rangle$}}

\newcommand{\green}{\textcolor{ForestGreen}}
\newcommand{\blue}{\textcolor{blue}}
\newcommand{\purple}{\textcolor{purple}}

\usepackage{algorithmic}            
\usepackage{algorithm}

\titleformat{\paragraph}
{\normalfont\normalsize\bfseries}{\theparagraph}{1em}{}
\titlespacing*{\paragraph}
{0pt}{3.25ex plus 1ex minus .2ex}{1.5ex plus .2ex}

\newcounter{savealgorithm}
\newenvironment{subalgorithms}
 {%
  \stepcounter{algorithm}%
  \edef\currentthealgorithm{\thealgorithm}%
  \setcounter{savealgorithm}{\value{algorithm}}%
  \setcounter{algorithm}{0}%
  \renewcommand{\thealgorithm}{\currentthealgorithm\alph{algorithm}}%
 }
 {%
  \setcounter{algorithm}{\value{savealgorithm}}%
 }
\begin{document}
\maketitle



%%%%%%%%%%%%%%%%%%%%%%%%%%%%%%%%%%%%%%%%%%%%%%%%%%%%%%%%%%%%%%%%%
\begin{abstract}
%%%%%%%%%%%%%%%%%%%%%%%%%%%%%%%%%%%%%%%%%%%%%%%%%%%%%%%%%%%%%%%%%


\end{abstract}

%%%%%



%%%%%%%%%%%%%%%%%%%%%%%%%%%%%%%%%%%%%%%%%%%%%%%%%%%%%%%%%%%%%%%%%
\section{Introduction}\label{sec:intro}
%%%%%%%%%%%%%%%%%%%%%%%%%%%%%%%%%%%%%%%%%%%%%%%%%%%%%%%%%%%%%%%%%

Some points to make

\begin{itemize}
    \item Other authors have modeled coevolution. \cite{gavrilets1998coevolutionary,ito2019joint}. 
    \item We introduce a stochastic models where we do not assume weak selection regimes or constant traits (e.g., effects of one species on the other). The traits in our models do evolve. In models like \cite{gavrilets1998coevolutionary}, the traits are constant. The presentation of this model would be Section 2 (Methods) of the manuscript.
    \item We are able to quantify the evolution on the level of cooperation (the $\theta$ values on our simulations). In \cite{ito2019joint} they also quantify cooperation. This could be a good paper to discuss.
    \item We use our model to examine how transitions between symbiotic scenarios (i.e., mutualism, parasitism, competition, predator-prey) occur. For example, we can evaluate the effect of disparities in mutation rates
          between hosts and endosymbionts on the emergence of a distinct symbiotic scenario than the initial symbiotic scenario. We can present generic simulations in a first subsection of the results section (which would be Subsection 3.1 of the manuscript).
    \item We can simplify the mutation mecanism to reproduce the four-state mutation framework in the paper of Mike \cite{lynch2023mutation}. We can then study the effect of differences in mutation rates. For example, we can examine scenarios when the evolution of one actor (i.e., host or endosymbiont) ``drives'' the evolution of the other. This is to be done.
         We could write the results of this part in a subsction of the results (e.g., Subsection 3.2)
\end{itemize}

%%%%%%%%%%%%%%%%%%%%%%%%%%%%%%%%%%%%%%%%%%%%%%%%%%%%%%%%%%%%%%%%
\section{Methods: the model}\label{sec:methods}
%%%%%%%%%%%%%%%%%%%%%%%%%%%%%%%%%%%%%%%%%%%%%%%%%%%%%%%%%%%%%%%%


 %%%%%%%%%%%%%%%%%%%%%%%%%%%%%%%%
\subsection*{Parameters and traits}
%%%%%%%%%%%%%%%%%%%%%%%%%%%%%%%%


 We summarize the model traits in Table \ref{table:traits}. 
 
 
\begin{table}[H]
\begin{centering}
\resizebox{\columnwidth}{!}{%
\begin{tabular}{lll}
{\bf Symbol} & {\bf Description} & {\bf Initial velues}  \\ \hline \hline
$s_{\text{host}}(t,i)$ & Baseline selection of host $i$ & $c =50$ \\ \hline
$h_{\text{host}}(t,i)$ & Effect of harvesting by invividual host $i$ on its selection &  0.25 \\ \hline
$e_{\text{host}}(t,i)$ & Effect of individual host $i$ on selection of its endosymbionts &  -0.75, -0.42, - 0.08, 0.25  \\ \hline
$s_{\text{sym}}(t,i,j)$ & Baseline selection of individual endosymbiont $j$ in its individual host $i$ & 1 \\ \hline
$h_{\text{sym}}(t,i,j)$ & Effect of harvesting by individual endosymbiont $j$ on its selection & 0.25 \\ \hline
$e_{\text{sym}}(t,i,j)$ & Effect of individual endosymbiont $j$ on selection of its host $i$ &  -0.75, -0.42, - 0.08, 0.25 \\ \hline 
$\theta_{\text{host}}$ & Mean of $h_{\text{host}}+e_{\text{sym}}$. Net symbiosis effect on host & -0.5, -0.17, - 0.17, 0.5 \\ \hline
$\theta_{\text{sym}}$ & Mean of $h_{\text{sym}} + e_{\text{host}}$. Net symbiosis effect on endosymbionts & -0.5, -0.17, - 0.17, 0.5 \\ \hline\hline
\end{tabular}
}
\end{centering}
\caption{Summary of traits in generation $t$.}
\label{table:traits}
\end{table}


We summarize the model parameters in Table \ref{table:parameters}. 
 
 
\begin{table}[H]
\begin{centering}
\begin{tabular}{lll}
{\bf Symbol} & {\bf Description} & {\bf Values}  \\ \hline \hline
  $N$ &  Number of hosts in each generation & 1000 \\ \hline 
  $M$ &  Number of free-living endosymbionts & 1000 \\ \hline 
  $c$ &  Maximum number of endosymbionts a single host can contain & 50 \\ \hline 
  $\gamma_{\text{sym}}$ &  Mutation variability for one endosymbiont & 0.05 \\ \hline 
  $\gamma_{\text{host}}$ &  Mutation variability for one host & 0.05 \\ \hline
  $\mu_{\text{sym}}$ &  Mutation distribution for endosymbionts & $\text{Uniform}\left( -\gamma_{\text{sym}},\gamma_{\text{sym}} \right)$ \\ \hline
  $\mu_{\text{host}}$ &  Mutation distribution for hosts & $\text{Uniform}\left( -\gamma_{\text{host}},\gamma_{\text{host}} \right)$ \\ \hline 
  $\epsilon_{\text{sym}}$ &  Mutation probability for one endosymbiont per generation & 0.001 or 0.00001 \\ \hline 
  $\epsilon_{\text{host}}$ &  Mutation probability for one host per generation & 0.001 or 0.00001 \\ \hline 
  $W_{\text{out}}$ &  Number of migrations between free-living and & 10 \\ 
  &  non free-living endosymbionts per generation &  \\ \hline 
  $W_{\text{in}}$ &  Number of migrations between endosymbionts & 10 \\
  & in different hosts per generation & \\ \hline 
  $D_{\text{host}}$ &  Death probability of a host per generation & 0.2 \\ \hline 
  $D_{\text{sym}}$ &  Death probability of a endosymbiont per generation  &  0.2 \\ \hline \hline
\end{tabular}
\end{centering}
\caption{Summary of parameters.}
\label{table:parameters}
\end{table}





%%%%%%%%%%%%%%%%%%%%%%%%%%%%%%%%
\subsection*{Fitness}
%%%%%%%%%%%%%%%%%%%%%%%%%%%%%%%%

\textbf{Free-living symbiont} $(g,0,j)$ has fitness:
\[
F(g,0,j) = s_{\text{sym}}(g,0,j)
\]
\textbf{Symbiont} $(g,i,j)$ (inside host $i \neq 0$) has fitness:
\[
F(g,i,j) = s_{\text{sym}}(g,i,j) + h_{\text{sym}}(g,i,j) + e_{\text{host}}(g,i)
\]
\textbf{host} $(g,i)$ has fitness:
\[
F(g,i) = s_{\text{host}}(g,i) + h_{\text{host}}(g,i) \cdot K_{g,i} + \sum_{j=1}^{K_{g,i}} e_{\text{sym}}(g,i,j)
\]
where $K_{g,i}$ is the number of alive symbionts hosted by host $(g,i)$.

%%%%%%%%%%%%%%%%%%%%%%%%%%%%%%%%
\subsection*{Dynamics per Generation}
%%%%%%%%%%%%%%%%%%%%%%%%%%%%%%%%

\begin{enumerate}
  \item \textbf{Death:}
  Each host dies with probability $D_{\text{host}}$; each symbiont with $D_{\text{sym}}$. Dead individuals are assigned trait $\delta$.

  \item \textbf{Reproduction:}
  \begin{itemize}
    \item Each dead host is replaced by a copy (with mutation) of a parent host from the previous generation, chosen with probability proportional to fitness. Dead hosts have fitness 1.
    \item Symbionts in living hosts are replaced from symbionts in the same location in the previous generation, also proportionally to fitness.
    \item Free-living symbionts are replaced from previous free-living symbionts.
    \item When a host reproduces, its symbionts are copied into both daughter hosts (for simplicity).
  \end{itemize}

  \item \textbf{Swapping:}
  \begin{itemize}
    \item Uniformly choose $W_{\text{out}}$ pairs of symbionts (one inside host, one free-living) to form set $\mathcal{W}_{\text{out}}$.
    \item Uniformly choose $W_{\text{in}}$ pairs of symbionts from different hosts to form set $\mathcal{W}_{\text{in}}$.
    \item Ensure disjointness: no symbiont appears in more than one pair.
    \item For each pair $((g,i,j), (g,r,l))$ in $\mathcal{W}_{\text{out}} \cup \mathcal{W}_{\text{in}}$, swap their traits.
  \end{itemize}

  \item \textbf{Mutation:}
After reproduction, each new individual receives mutated traits according to either $\mu_{\text{host}}$ or $\mu_{\text{sym}}$, depending on its type. 

To be concrete, assume that $\mu_{\text{host}}$ and $\mu_{\text{sym}}$ are independent, and that
\[
\mu_{\text{host}} \sim (1 - \epsilon_{\text{host}})\, \delta_0 + \epsilon_{\text{host}}\, \delta_{(U_1, U_2, U_3)},
\]
where $(U_1, U_2, U_3)$ is a random vector uniformly distributed on $[-\gamma_{\text{host}}, \gamma_{\text{host}}]^3$. Similarly, let
\[
\mu_{\text{sym}} \sim (1 - \epsilon_{\text{sym}})\, \delta_0 + \epsilon_{\text{sym}}\, \delta_{(V_1, V_2, V_3)},
\]
where $(V_1, V_2, V_3)$ is also uniformly distributed on $[-\gamma_{\text{sym}}, \gamma_{\text{sym}}]^3$.

If a newly born host has a mother with traits 
\[
(s_{\text{host}}(g,i),\ h_{\text{host}}(g,i),\ e_{\text{host}}(g,i)),
\]
then with probability $1 - \epsilon_{\text{host}}$ the offspring inherits the exact same traits. Otherwise, it mutates and its traits become:
\[
\left(
s_{\text{host}}(g,i) + U_1,\ 
h_{\text{host}}(g,i) + U_2,\ 
e_{\text{host}}(g,i) + U_3
\right).
\]
Similarly, a symbiont inherits the exact traits of its mother with probability $1 - \epsilon_{\text{sym}}$, and otherwise mutates according to:
\[
\left(
s_{\text{sym}}(g,i) + V_1,\ 
h_{\text{sym}}(g,i) + V_2,\ 
e_{\text{sym}}(g,i) + V_3
\right),
\]
where $(V_1, V_2, V_3)$ is drawn uniformly from $[-\gamma_{\text{sym}}, \gamma_{\text{sym}}]^3$ as defined above, and $(s_{\text{sym}}, h_{\text{sym}}, e_{\text{sym}})$ denote the corresponding traits of the symbiont mother. 



\end{enumerate}

%%%%%%%%%%%%%%%%%%%%%%%%%%%%%%%%%%%%%%%%%%%%%%%%%%%%%%%%%%%%%%%%
\section{Results}\label{sec:results}
%%%%%%%%%%%%%%%%%%%%%%%%%%%%%%%%%%%%%%%%%%%%%%%%%%%%%%%%%%%%%%%%


%%%%%%%%%%%%%%%%%%%%%%%%%%%%%%%%%%%%%%%%%%%%%%%%%%%%%%%%%%%%%%%%
\subsection{Transitions across four scenarios of symbiosis}\label{sec:four_scenarios}
%%%%%%%%%%%%%%%%%%%%%%%%%%%%%%%%%%%%%%%%%%%%%%%%%%%%%%%%%%%%%%%%

\begin{itemize}
    \item In the following simulations we combine the trait initial conditions in Table \ref{table:traits} and the paremeters in Table \ref{table:parameters}
          to generate $N_{sim} = 3$ simulations for each combination during $N_g = 10000$ generations.
    \item We keep track of trajectories of the net effect of symbiosis with the $\theta$ values (panels (a) in Figures 1, 3, 4).
    \item In Figure 2 we keep track of the populations of host and symbionts for initial conditions of the competition scenario. We see that curves move towards
          mutualism when the non free-living populations are large and these curves remain in competition when the symbionts populations collapse.
    \item We keep track of the final scenario (at the last generation) for each set of simulations. We plot the distributions of the final scenarios (over three simulations) for different
          conbinations of initial conditions on panels (b) in Figures 1, 3, 4.
    \item We keep track of the number of generations in each scenario over different simulations. We plot the distributions of the spent time in all the scenarios for different
          conbinations of initial conditions on panels (c) in Figures 1, 3, 4.
    \item General pattern: convergence to mutualism. How this convergence takes place depends on relative differences in mutation rates of host and symbionts and initial scenario.            
\end{itemize}

%%%%%%%%%%%%%%%%%%%%%%%
\begin{figure}[H]
    \centering
    \subfloat[][Evolution of trajectories]{\includegraphics[scale=0.6]{../Figures/theta_trajectories_epsilon_host_0.001_epsilon_sym_0.001.png}}\\
    \subfloat[][Final scenario distributions]{\includegraphics[scale=0.42]{../Figures/final_scenario_epsilon_cell_0.001_epislon_sym_0.001.pdf}}
    \subfloat[][Spent time in each scenario]{\includegraphics[scale=0.42]{../Figures/distribution_duration_each_scenario_epsilon_cell_0.001_epislon_sym_0.001.pdf}}
    \caption{Coevolution of the parameters $\theta_{\text{host}}$, $\theta_{\text{sym}}$ where $\epsilon_{\text{host}} = 0.001$ and $\epsilon_{\text{sym}} = 0.001$. }
    \label{fig:absorption_graph}
\end{figure}
%%%%%%%%%%%%%%%%%%%%%%%

\begin{figure}[H]
    \centering
    \subfloat[][]{\includegraphics[scale=0.5]{../Figures/alive_sym_in_epsilon_host_0.001_epsilon_sym_0.001_theta_host_-0.5_theta_sym_-0.5.png}}
    \subfloat[][]{\includegraphics[scale=0.5]{../Figures/alive_sym_in_epsilon_host_0.001_epsilon_sym_0.001_theta_host_-0.17_theta_sym_-0.17.png}}\\
    \subfloat[][]{\includegraphics[scale=0.5]{../Figures/thetas_zoom_in_epsilon_host_0.001_epsilon_sym_0.001_theta_host_-0.5_theta_sym_-0.5.png}}
    \subfloat[][]{\includegraphics[scale=0.5]{../Figures/thetas_zoom_in_epsilon_host_0.001_epsilon_sym_0.001_theta_host_-0.17_theta_sym_-0.17.png}}
    \caption{Number of symbionts and $\theta$-trajectories for simulations that starting in two {\bf competition} scenarios. Here $\epsilon_{\text{host}} = 0.001$ and $\epsilon_{\text{sym}} = 0.001$. }
    \label{fig:absorption_graph}
\end{figure}
%%%%%%%%%%%%%%%%%%%%%%%



\begin{figure}[H]
    \centering
    \subfloat[][Evolution of trajectories]{\includegraphics[scale=0.6]{../Figures/theta_trajectories_epsilon_host_0.001_epsilon_sym_1e-05.png}}\\
    \subfloat[][Final scenario distributions]{\includegraphics[scale=0.42]{../Figures/final_scenario_epsilon_cell_0.001_epislon_sym_1e-05.pdf}}
    \subfloat[][Spent time in each scenario]{\includegraphics[scale=0.42]{../Figures/distribution_duration_each_scenario_epsilon_cell_0.001_epislon_sym_1e-05.pdf}}
    \caption{Coevolution of the parameters $\theta_{\text{host}}$, $\theta_{\text{sym}}$ where $\epsilon_{\text{host}} = 0.001$ and $\epsilon_{\text{sym}} = 0.00001$. }
    \label{fig:absorption_graph}
\end{figure}
%%%%%%%%%%%%%%%%%%%%%%%

\begin{figure}[H]
    \centering
    \subfloat[][Evolution of trajectories]{\includegraphics[scale=0.6]{../Figures/theta_trajectories_epsilon_host_1e-05_epsilon_sym_0.001.png}}\\
    \subfloat[][Final scenario distributions]{\includegraphics[scale=0.42]{../Figures/final_scenario_epsilon_cell_1e-05_epislon_sym_0.001.pdf}}
    \subfloat[][Spent time in each scenario]{\includegraphics[scale=0.42]{../Figures/distribution_duration_each_scenario_epsilon_cell_1e-05_epislon_sym_0.001.pdf}}
    \caption{Coevolution of the parameters $\theta_{\text{host}}$, $\theta_{\text{sym}}$ where $\epsilon_{\text{host}} = 0.00001$ and $\epsilon_{\text{sym}} = 0.001$. }
    \label{fig:absorption_graph}
\end{figure}
%%%%%%%%%%%%%%%%%%%%%%%

\begin{figure}[H]
    \centering
    \subfloat[][Evolution of trajectories]{\includegraphics[scale=0.6]{../Figures/theta_trajectories_epsilon_host_1e-05_epsilon_sym_1e-05.png}}\\
    \subfloat[][Final scenario distributions]{\includegraphics[scale=0.42]{../Figures/final_scenario_epsilon_cell_1e-05_epislon_sym_1e-05.pdf}}
    \subfloat[][Spent time in each scenario]{\includegraphics[scale=0.42]{../Figures/distribution_duration_each_scenario_epsilon_cell_1e-05_epislon_sym_1e-05.pdf}}
    \caption{Coevolution of the parameters $\theta_{\text{host}}$, $\theta_{\text{sym}}$ where $\epsilon_{\text{host}} = 0.00001$ and $\epsilon_{\text{sym}} = 0.00001$. }
    \label{fig:absorption_graph}
\end{figure}
%%%%%%%%%%%%%%%%%%%%%%%


 
 
%%%%%%%%%%%%%%%%%%%%%%%%%%%%%%%%%%%%%%%%%%%%%%%%%%%%%%%%%%%%%%%%
\subsection{Competition in a simple mutatation regime}\label{sec:subResult}
%%%%%%%%%%%%%%%%%%%%%%%%%%%%%%%%%%%%%%%%%%%%%%%%%%%%%%%%%%%%%%%%


%%%%%%%%%%%%%%%%%%%%%%%%%%%%%%%%%%%%%%%%%%%%%%%%%%%%%%%%%%%%%%%%
\section{Conclusions and Discussion}\label{sec:discussion}
%%%%%%%%%%%%%%%%%%%%%%%%%%%%%%%%%%%%%%%%%%%%%%%%%%%%%%%%%%%%%%%%
%%%%%

\appendix

%%%%%%%%%%%%%%%%%%%%%%%%%%%%%%%%%%%%%%%%%%%%%%%%%%%%%%%%%%%%%%%%
\section{Appendix}\label{appendix}
%%%%%%%%%%%%%%%%%%%%%%%%%%%%%%%%%%%%%%%%%%%%%%%%%%%%%%%%%%%%%%%%


\begin{algorithm}
\caption{Symbiosis model}
\label{alg:model}
\begin{algorithmic}[1]
 \renewcommand{\algorithmicrequire}{Input:}
 \renewcommand{\algorithmicensure}{Output:}
 \REQUIRE Inputs. \medskip
 \ENSURE Outputs.
\medskip \medskip
	\STATE Step 1
	\STATE Spep 2
	\STATE Step 3
\end{algorithmic}
\end{algorithm}


%%%%%%%%%%%%%%%%%%%%%%%%%%%%%%%%%%%%%%%%%%%%%%%%%%%%%%%%%%%%%%%%

\section*{Data and code availability}

%We have posted t
The code that yields the numerical results in Figures \ref{fig:} and \ref{fig:} of Section \ref{sec:} are available at \url{}.


%%%%%%%%%%%%%%%%%%%%%%%%%%%%%%%%%%%%%%%%%%%%%%%%%%%%%%%%%%%%%%%%
\section*{Acknowledgements}
%%%%%%%%%%%%%%%%%%%%%%%%%%%%%%%%%%%%%%%%%%%%%%%%%%%%%%%%%%%%%%%%

%This work was supported by 
We thank ....
%%%%%%%

\bibliography{biblio}
\bibliographystyle{plain}


\end{document}



